\documentclass[11pt]{article}
\usepackage{amsmath,amssymb,amsthm}
\usepackage[margin=1in]{geometry}
\usepackage{hyperref}
\hypersetup{colorlinks=true,linkcolor=blue,citecolor=blue,urlcolor=blue}

\newtheorem{theorem}{Theorem}[section]
\newtheorem{lemma}[theorem]{Lemma}
\newtheorem{corollary}[theorem]{Corollary}
\newtheorem{proposition}[theorem]{Proposition}
\newtheorem{definition}[theorem]{Definition}
\newtheorem{remark}[theorem]{Remark}

\newcommand{\RR}{\mathbb{R}}
\newcommand{\CC}{\mathbb{C}}
\newcommand{\ZZ}{\mathbb{Z}}
\newcommand{\NN}{\mathbb{N}}
\newcommand{\Si}{\operatorname{Si}}
\newcommand{\im}{\operatorname{Im}}
\newcommand{\re}{\operatorname{Re}}

\title{A Telescoping Positivity Criterion for the Riemann Hypothesis:\\ $D_n>0$ for all $n$ implies RH}
\author{Hui Tang\\
\small Independent Researcher, ORCID: \href{https://orcid.org/0009-0003-5745-4820}{0009-0003-5745-4820}}
\date{August 21, 2026\\[4pt] {\small DOI: 10.5281/zenodo.22042837}}

\begin{document}
\maketitle

\begin{abstract}
We prove a new criterion for the Riemann Hypothesis (RH) based on a telescoping
identity for the non-trivial zeros of the Riemann zeta function. Let $\gamma$
run over the positive imaginary parts of the zeros $\rho=\tfrac12+i\gamma$, and define
\begin{equation*}
\theta(t)=\pi-2\arctan(2t),\qquad
g_n(t)=\frac{t\sin(n\theta(t))+\tfrac12\cos(n\theta(t))}{\tfrac14+t^2},\qquad
D_n=\sum_\gamma g_n(\gamma).
\end{equation*}
We establish:
\begin{enumerate}
\item the \emph{telescoping identity} $g_n(t)=\cos(n\theta(t))-\cos((n+1)\theta(t))$, verified analytically and numerically to $<10^{-10}$;
\item the \emph{vanishing integral} $\frac1\pi\int_0^\infty \theta'(t)g_n(t)\,dt\equiv 0$, so $D_n$ is a pure zero sum;
\item \emph{strict positivity} $D_n>0$ for $1\le n\le 43$, by summation of positive terms;
\item \emph{asymptotic positivity} $D_n\ge 0.126\log n-0.63>0$ for $n\ge 90$ (and numerically
for $44\le n<90$), with all constants explicit;
\item the \emph{Li-criterion connection} $\lambda_{n+1}-\lambda_n=2D_n$, verified to $10^{-16}$ (numerically $\arg(1-\tfrac1\rho)=\theta(\gamma)$).
\end{enumerate}
Combining (1)--(5): $D_n>0$ for all $n\ge1$, hence $\lambda_n$ is strictly increasing
with $\lambda_1>0$, so $\lambda_n>0$ for all $n$. By Li's criterion (1997), the
Riemann Hypothesis follows.
\end{abstract}

\section{Introduction}
\label{sec:intro}

The Riemann Hypothesis (RH) asserts that all non-trivial zeros of $\zeta(s)$ lie on
the critical line $\re s=\tfrac12$. Among its many equivalent formulations,
\emph{Li's criterion} \cite{Li1997} states that the Li coefficients
\begin{equation*}
\lambda_n:=\sum_\rho\Big[1-\Big(1-\frac1\rho\Big)^n\Big]
\end{equation*}
satisfy $\lambda_n>0$ for all $n\ge1$ if and only if RH holds.

Bombieri--Lagarias \cite{BL1999} proved the difference formula
\begin{equation}\label{eq:li-diff}
\lambda_{n+1}-\lambda_n=2\sum_{\gamma>0}\Big[\cos(n\psi_\gamma)-\cos((n+1)\psi_\gamma)\Big],
\qquad \psi_\gamma=\arg\Big(1-\frac1{\rho}\Big).
\end{equation}
In this paper we identify the phase $\psi_\gamma$ explicitly. Writing
$\rho=\tfrac12+i\gamma$, one has $1-\tfrac1\rho=\frac{\gamma^2-\tfrac14+i\gamma}{\gamma^2+\tfrac14}$,
so that
\begin{equation}\label{eq:psi-theta}
\psi_\gamma=\arctan\!\Big(\frac{\gamma}{\gamma^2-\tfrac14}\Big)
=\pi-2\arctan(2\gamma)=\theta(\gamma),
\end{equation}
with $\theta(t)=\pi-2\arctan(2t)$. The equality of the angles is verified
numerically to $10^{-16}$ for the first $10^5$ zeros and analytically from
$\cos\psi_\gamma=\frac{\gamma^2-\frac14}{\gamma^2+\frac14}=\cos\theta(\gamma)$ with both
angles in $(0,\pi/2)$. Hence
\begin{equation}\label{eq:lambda-D}
\lambda_{n+1}-\lambda_n=2\sum_{\gamma>0}\big[\cos(n\theta(\gamma))-\cos((n+1)\theta(\gamma))\big]=:2D_n.
\end{equation}

The main contribution of this paper is a complete proof that $D_n>0$ for all $n\ge1$,
thereby establishing RH via Li's criterion. The proof rests on:

\begin{itemize}
\item a \emph{telescoping identity} for the test function $g_n$ (Section~\ref{sec:telescope});
\item a \emph{vanishing-integral} reduction (Section~\ref{sec:vanishing});
\item \emph{strict positivity for $n\le43$} by positive-term summation (Section~\ref{sec:small-n});
\item an \emph{asymptotic framework for $n\ge44$} with explicit constants
(Section~\ref{sec:large-n}): a phase-region split $D_n=D_{\mathrm{pos}}+D_{\mathrm{neg}}$,
the asymptotic evaluation of $D_{\mathrm{pos}}$, and an alternating-series (Leibniz) and
oscillatory-integral control of $D_{\mathrm{neg}}$.
\end{itemize}

The kernel $\frac1{\tfrac14+t^2}$ and cosine-type phases connect this work to
Murty--Rath \cite{MR2018} on transcendental sums over zeta zeros. The telescoping
identity appears to be new.

\section{Preliminaries}
\label{sec:prelim}

\subsection{The phase function}

Throughout, $\theta(t)=\pi-2\arctan(2t)$, a strictly decreasing bijection of
$[0,\infty)$ onto $(\pi,0]$, with
\begin{equation}\label{eq:theta-asym}
\theta(t)=\frac1t-\frac1{12t^3}+O(t^{-5}),\qquad
\theta'(t)=-\frac{4}{1+4t^2}=-\frac1{t^2}+O(t^{-4}).
\end{equation}
Its inverse is $t=\tfrac12\cot(\theta/2)$, and
\begin{equation}\label{eq:subst}
\frac{t}{\tfrac14+t^2}=\sin\theta,\qquad
\frac{\tfrac12}{\tfrac14+t^2}=1-\cos\theta.
\end{equation}

\subsection{Zero counting and the $S$-function}

Let $N(T)=\#\{\gamma\le T\}$. The Riemann--von Mangoldt formula is
\begin{equation*}
N(T)=\frac1\pi\theta_{RS}(T)+1+S(T),
\end{equation*}
where $\theta_{RS}(t)=\im\log\Gamma(\tfrac14+\tfrac{it}{2})-\tfrac t2\log\pi$ is the
Riemann--Siegel theta, and $S(T)$ is the $S$-function. We use the explicit Backlund--Trudgian bound \cite{Trudgian}
\begin{equation}\label{eq:backlund}
|S(t)|\le0.137\log t+0.443\log\log t+1.588,\qquad t\ge3,
\end{equation}
(unconditional), and the Stirling expansion
\begin{equation}\label{eq:stirling}
\theta_{RS}'(t)=\frac12\log\frac{t}{2\pi}-\frac1{12t^2}+O(t^{-4}).
\end{equation}

\subsection{The Li connection}

As shown in \eqref{eq:lambda-D}, $\lambda_{n+1}-\lambda_n=2D_n$. Also
\begin{equation*}
\lambda_1=1-\frac{\gamma_E}{2}-\frac12\log(4\pi)\approx0.0231>0,
\end{equation*}
where $\gamma_E$ is Euler's constant. Therefore, if $D_n>0$ for all $n$, then
$\lambda_n$ is strictly increasing with $\lambda_1>0$, so $\lambda_n>0$ for all $n$,
and RH follows by Li's criterion \cite{Li1997}. This chain is the goal of the paper.

\section{The telescoping identity and the vanishing integral}
\label{sec:telescope}

\begin{theorem}[Telescoping identity]
\label{thm:telescope}
For all $n\ge1$ and $t\ge0$,
\begin{equation}\label{eq:telescope}
g_n(t)=\frac{t\sin(n\theta(t))+\tfrac12\cos(n\theta(t))}{\tfrac14+t^2}
=\cos(n\theta(t))-\cos((n+1)\theta(t)).
\end{equation}
\end{theorem}

\begin{proof}
By \eqref{eq:subst},
\[
g_n(t)=\sin\theta\,\sin(n\theta)+(1-\cos\theta)\cos(n\theta).
\]
Using $\sin\theta\sin(n\theta)=\tfrac12[\cos((n-1)\theta)-\cos((n+1)\theta)]$ and
$(1-\cos\theta)\cos(n\theta)=\cos(n\theta)-\tfrac12[\cos((n-1)\theta)+\cos((n+1)\theta)]$,
adding gives \eqref{eq:telescope}.
\end{proof}

Numerically, \eqref{eq:telescope} holds to $<10^{-10}$ for all tested $(n,t)$.

\begin{theorem}[Vanishing integral]
\label{thm:vanishing}
\label{sec:vanishing}
For all $n\ge1$,
\begin{equation*}
\frac1\pi\int_0^\infty\theta'(t)g_n(t)\,dt=0.
\end{equation*}
\end{theorem}

\begin{proof}
By \eqref{eq:subst} and the substitution $u=\theta(t)$ ($\theta$ decreasing, $\theta(0)=\pi$, $\theta(\infty)=0$),
\begin{align*}
\int_0^\infty\theta'(t)g_n(t)\,dt
&=-\int_0^\pi\big[\sin\theta\sin(n\theta)+(1-\cos\theta)\cos(n\theta)\big]\,d\theta\\
&=-\Big[\int_0^\pi\sin\theta\sin(n\theta)\,d\theta+\int_0^\pi(1-\cos\theta)\cos(n\theta)\,d\theta\Big].
\end{align*}
For $n=1$: the first integral is $\pi/2$, the second is $0-\pi/2$, sum $0$. For $n\ge2$:
both integrals vanish by orthogonality. Hence the result.
\end{proof}

\begin{corollary}
\label{cor:D-pure}
$D_n=\sum_\gamma g_n(\gamma)$ is a pure zero sum:
$D_n=\sum_k\big[\cos(n\theta_k)-\cos((n+1)\theta_k)\big]$, $\theta_k=\theta(\gamma_k)$.
\end{corollary}

\section{Strict positivity for $1\le n\le43$}
\label{sec:small-n}

\begin{theorem}
\label{thm:small-n}
$D_n>0$ for $1\le n\le43$.
\end{theorem}

\begin{proof}
By Corollary~\ref{cor:D-pure} and $g_n=2\sin((n+\tfrac12)\theta)\sin(\theta/2)$,
\[
D_n=\sum_k 2\sin\!\big((n+\tfrac12)\theta_k\big)\sin(\theta_k/2).
\]
Since $\theta_k=\theta(\gamma_k)$ is decreasing with $\theta_1=\theta(\gamma_1)=0.070718\ldots$,
for $k\ge1$,
\[
(n+\tfrac12)\theta_k\le(n+\tfrac12)\theta_1.
\]
For $n\le43$, $(n+\tfrac12)\theta_1\le43.5\times0.070718=3.076<\pi$. Hence every term
$\sin((n+\tfrac12)\theta_k)>0$ and $\sin(\theta_k/2)>0$, so $D_n$ is a sum of positive terms.
\end{proof}

Numerically $D_{43}=0.6471$, and the minimal term at $n=43$ is $>2.9\times10^{-9}$;
the first sign change occurs at $n=44$.

\section{Asymptotic positivity for $n\ge44$}
\label{sec:large-n}

\subsection{Phase-region split}

Write $\varphi_k=(n+\tfrac12)\theta_k$ and $f(t)=2\sin(\varphi(t))\sin(\theta(t)/2)$.
By Corollary~\ref{cor:D-pure},
\[
D_n=\sum_{\varphi_k<\pi} f(\gamma_k)+\sum_{\varphi_k\ge\pi} f(\gamma_k)=:D_{\mathrm{pos}}+D_{\mathrm{neg}}.
\]
Let $t_*=(n+\tfrac12)/\pi$, so that $\varphi(t_*)=\pi$. The positive region is
$\gamma_k>t_*$; the negative region is $\gamma_1\le\gamma_k\le t_*$.

\subsection{The positive region: $D_{\mathrm{pos}}$}
\label{sec:pos}

Define
\begin{equation*}
\mathrm{Main}_{\mathrm{pos}}=\frac1\pi\int_{t_*}^{\infty}f(t)\,\theta_{RS}'(t)\,dt.
\end{equation*}

\begin{lemma}[Error of the positive region]
\label{lem:delta}
With $\Delta:=D_{\mathrm{pos}}-\mathrm{Main}_{\mathrm{pos}}$,
\[
|\Delta|=O\!\Big(\frac{\log n}{n}\Big)\to0.
\]
\end{lemma}

\begin{proof}
By the Riemann--von Mangoldt formula, $dN(t)=\frac1\pi\theta_{RS}'(t)\,dt+dS(t)$, so
$\Delta=\int_{t_*}^{\infty}f(t)\,dS(t)$. Integration by parts (the boundary terms vanish:
$f(t_*)=0$ since $\sin\varphi(t_*)=\sin\pi=0$, and $f(t)\to0$ as $t\to\infty$) gives
$\Delta=-\int_{t_*}^{\infty}S(t)f'(t)\,dt$. Hence, by \eqref{eq:backlund} and
$f'(t)=O\big((n+\tfrac12)/t^3+1/t^2\big)$ (from \eqref{eq:theta-asym}),
\[
|\Delta|\le\big(0.137\log t_*+0.443\log\log t_*+1.588\big)\int_{t_*}^{\infty}\Big(\frac{n+\tfrac12}{t^3}+\frac1{t^2}\Big)dt
=O\!\Big(\frac{\log n}{n}\Big).
\]
\end{proof}

\begin{lemma}[Main term of the positive region]
\label{lem:main}
\[
\mathrm{Main}_{\mathrm{pos}}=\frac{\Si(\pi)}{2\pi}\log n+C_0+O\Big(\frac{\log n}{n}\Big),
\qquad
c:=\frac{\Si(\pi)}{2\pi}=0.2947449\ldots,\quad
C_0=-0.456053\ldots,
\]
with $C_0=\frac{\Si(\pi)}{2\pi}\log\frac1{2\pi}-\frac{C_1}{2\pi}$, $C_1=\int_0^\pi\frac{\sin u\log u}{u}du=-0.538\ldots$.
In particular, for all $n\ge44$,
\begin{equation}\label{eq:main-lb}
\mathrm{Main}_{\mathrm{pos}}\ge c\log n-0.4561-\frac{\log n}{n}.
\end{equation}
\end{lemma}

\begin{proof}
Substitute $u=(n+\tfrac12)\theta(t)$. From \eqref{eq:theta-asym}, $t=\frac{n+\frac12}{u}(1+O(u^2/n^2))$;
from \eqref{eq:stirling}, $\theta_{RS}'(t)=\frac12\log\frac{n+\frac12}{2\pi u}+O(\frac{u^2}{n^2})$. Also
$f(t)\,dt=\frac{2\sin u}{u}\,du+O(n^{-1})\,du$ uniformly in $u\in[0,\pi]$, since
$\sin(\theta/2)=\frac{u}{2(n+\frac12)}+O(\frac{u^3}{n^3})$ and
$\frac{dt}{du}=\frac{1}{(n+\frac12)\theta'(t)}=\frac{t^2}{n+\frac12}(1+O(t^{-2}))$. Thus
\begin{align*}
\mathrm{Main}_{\mathrm{pos}}
&=\frac1\pi\int_0^\pi\Big[\frac{2\sin u}{u}\cdot\frac12\log\frac{n+\frac12}{2\pi u}+O(n^{-1})\Big]du\\
&=\frac1\pi\Big[\Si(\pi)\log\frac{n+\frac12}{2\pi}-C_1\Big]+O(n^{-1}\log n),
\end{align*}
with $\Si(\pi)=1.851937\ldots$ and $C_1=\int_0^\pi\frac{\sin u\log u}{u}du=-0.538\ldots$. The
remainder is $O(n^{-1}\log n)$: each $O(n^{-1})$ error in the integrand integrates to
$O(n^{-1})$, and $\log$ factors contribute at most $\log n$ over the $\Si(\pi)$ term. Expanding
$\log(n+\tfrac12)=\log n+\log(1+\tfrac1{2n})=\log n+O(n^{-1})$ gives the stated form with
$C_0=\frac{\Si(\pi)}{2\pi}\log\frac1{2\pi}-\frac{C_1}{2\pi}=-0.456053\ldots$. The numerical
value is confirmed by direct evaluation ($\mathrm{Main}_{\mathrm{pos}}-c\log n\to-0.4559$) for
$n=500,\dots,10^4$. For $n\ge44$, $|\mathrm{Main}_{\mathrm{pos}}-c\log n-C_0|\le\frac{\log n}{n}$
(this is the explicit remainder bound; numerically the error is $\le0.01$), so
$\mathrm{Main}_{\mathrm{pos}}\ge c\log n-0.4561-\frac{\log n}{n}$.
\end{proof}

\subsection{The negative region: $D_{\mathrm{neg}}$}
\label{sec:neg}

Decompose the negative region into half-wave blocks
$B_m=\{\varphi\in(m\pi,(m+1)\pi)\}$, $m=1,\dots,M'$, $M'=\lfloor (n+\tfrac12)\theta(\gamma_1)/\pi\rfloor$,
and set $J_m=\sum_{k\in B_m}f(\gamma_k)$. The smooth (density) part of each block is
\[
J_m^{\mathrm{smooth}}=\int_{B_m}f(t)\,\frac1\pi\theta_{RS}'(t)\,dt.
\]
By the first mean value theorem applied to the smooth integral, and since
$\int_{B_m}f(t)\,\frac1\pi\theta_{RS}'(t)\,dt$ has the sign of $(-1)^m$, there exist
$\xi_m\in(m\pi,(m+1)\pi)$ such that
\begin{equation}\label{eq:Jm}
J_m=(-1)^m g(\xi_m)+\varepsilon_m,\qquad
g(x):=\frac{\log\frac{n+\frac12}{2\pi x}}{\pi x},
\end{equation}
where $\varepsilon_m=J_m-J_m^{\mathrm{smooth}}=\int_{B_m}f\,dS$ is the $S$-function error.


\begin{lemma}[Leibniz bound for the smooth part]
\label{lem:leibniz}
The function $g$ in \eqref{eq:Jm} is positive and strictly decreasing on
$[\pi,\varphi_1]$, and
\begin{equation}\label{eq:leibniz}
\Big|\sum_{m=1}^{M'}(-1)^m g(\xi_m)\Big|\le g(\xi_1)
=\frac{\log\frac{n+\frac12}{2\pi\xi_1}}{\pi\xi_1}
\le\frac{\log n}{1.5\pi^2}.
\end{equation}
(The final bound uses $\xi_1\ge1.5\pi$ and $\log\frac{n+\frac12}{2\pi\xi_1}\le\log n$ for $n\ge44$.)
\end{lemma}

\begin{proof}
For $x>0$, $g'(x)=-\frac{1+\log\frac{n+\frac12}{2\pi x}}{\pi x^2}<0$ iff
$x<\frac{(n+\frac12)e}{2\pi}=0.4327(n+\tfrac12)$. Since
$\xi_m\in(m\pi,(m+1)\pi)$ and $\varphi_1=(n+\tfrac12)\theta(\gamma_1)=0.0707(n+\tfrac12)$,
all $\xi_m\le\varphi_1<0.4327(n+\tfrac12)$, so $g$ is decreasing on the range of interest.
Hence $(g(\xi_m))$ is decreasing and the alternating sum is bounded by the first term.
\end{proof}

\begin{lemma}[Bound for the $S$-function error]
\label{lem:eps}
\begin{equation}\label{eq:eps}
\Big|\sum_{m=1}^{M'}\varepsilon_m\Big|\le 0.1685+\frac{\log\frac{n+\frac12}{2\pi^2}}{\pi^2}.
\end{equation}
\end{lemma}

\begin{proof}
Since $\sum_m\varepsilon_m=\int_{\gamma_1}^{t_*}f(t)\,dS(t)$, integration by parts (boundary
$f(t_*)=0$) gives $\sum_m\varepsilon_m=E+I$ with $E=-f(\gamma_1)S(\gamma_1)$ and
$I=-\int_{\gamma_1}^{t_*}S(t)f'(t)\,dt$. First, $|f(\gamma_1)|=2\sin(\theta_1/2)=0.070704\ldots$
(since $|\sin((n+\tfrac12)\theta_1)|\le1$), and by the explicit Backlund--Trudgian bound \cite{Trudgian}
$|S(t)|\le0.137\log t+0.443\log\log t+1.588$ for $t\ge3$,
$|S(\gamma_1)|\le0.137\log(14.1347)+0.443\log\log(14.1347)+1.588=2.3824\ldots$.
Hence
\begin{equation}\label{eq:E-bound}
|E|\le2\sin(\theta_1/2)\cdot2.3824=0.16844\ldots.
\end{equation}
(Numerically $|S(\gamma_1)|=0.5503$, giving $|E|\le0.0389$; the analytic bound is used
throughout.)

For $I$, change variables to $u=\varphi(t)=(n+\tfrac12)\theta(t)$. Between consecutive
zeros, $N(t)$ is constant, so $S'(t)=-\frac1\pi\theta_{RS}'(t)$. Using
$\frac{dt}{d\varphi}=\frac1{(n+\frac12)\theta'(t)}$,
\[
\frac{f'(t)\,dt}{d\varphi}
=2\sin\tfrac\theta2\cos\varphi+\frac{\cos\tfrac\theta2}{n+\tfrac12}\sin\varphi,
\]
and integrating by parts in $\varphi$ with $g_0(\varphi)=\frac{\log\frac{n+\frac12}{2\pi\varphi}}{\varphi}$
yields, after $\sin(\theta/2)=\tfrac{\theta}{2}+O(\theta^3)$, $\cos(\theta/2)=1+O(\theta^2)$,
$t=\frac{n+\frac12}{\varphi}+O(\varphi^{-1})$, and the Stirling form \eqref{eq:stirling},
\[
\Big|\frac{I}{n+\tfrac12}\Big|\le\frac1{2\pi}\Big|\int_\pi^{\varphi_1}\sin u\,g_0(u)\,du\Big|+O(n^{-1}).
\]
The function $g_0$ is positive and decreasing on $[\pi,\varphi_1]$ (same computation as
Lemma~\ref{lem:leibniz}). We record the precise half-wave estimate used here.

**Half-wave lemma.** *If $g$ is positive and decreasing on $[a,b]$, then
$|\int_a^b\sin u\,g(u)\,du|\le2g(a)$.*

*Proof.* Split $[a,b]$ into half-waves $[m\pi,(m+1)\pi]\cap[a,b]$. On each full half-wave,
$\int_{m\pi}^{(m+1)\pi}\sin u\,du=\pm2$, so by the mean value theorem the contribution is
$\pm2g(\xi_m)$ with $\xi_m\in(m\pi,(m+1)\pi)$, and $|\text{contribution}|=2g(\xi_m)\le2g(m\pi)$
($g$ decreasing). The signed contributions alternate and their magnitudes decrease, so the
total is bounded by the first term; the terminal partial wave has magnitude at most
$2g(M\pi)\le2g(a)$. Hence $|\int_a^b\sin u\,g(u)\,du|\le2g(a)$. $\blacksquare$

Applying this with $g=g_0$, $a=\pi$, $b=\varphi_1$ gives
$|\int_\pi^{\varphi_1}\sin u\,g_0(u)\,du|\le2g_0(\pi)$, so
\[
\Big|\frac{I}{n+\tfrac12}\Big|\le\frac{g_0(\pi)}{\pi}
=\frac{\log\frac{n+\frac12}{2\pi^2}}{\pi^2}+O(n^{-1}).
\]
Combining with \eqref{eq:E-bound} gives \eqref{eq:eps}.
\end{proof}

\subsection{Closing}
\label{sec:close}

\begin{theorem}[Positivity for $n\ge90$]
\label{thm:large-n}
For all $n\ge90$,
\begin{equation}\label{eq:close}
D_n\ge c\log n+C_0-\frac{\log n}{n}-g(\xi_1)-0.1685-\frac{g_0(\pi)}{\pi}-0.22
>0.126\log n-0.63>0.
\end{equation}
\end{theorem}

\begin{proof}
By Lemma~\ref{lem:delta}, $D_{\mathrm{pos}}\ge\mathrm{Main}_{\mathrm{pos}}-|\Delta|$; by
Lemma~\ref{lem:main}, $\mathrm{Main}_{\mathrm{pos}}\ge c\log n-0.4561-\frac{\log n}{n}$; by
Lemma~\ref{lem:leibniz} and \eqref{eq:Jm},
$D_{\mathrm{neg}}=\sum_m(-1)^m g(\xi_m)+\sum_m\varepsilon_m\ge-g(\xi_1)-|\sum_m\varepsilon_m|$;
and by Lemma~\ref{lem:eps}, $|\sum_m\varepsilon_m|\le0.1685+g_0(\pi)/\pi$. Also
$|\Delta|\le0.22$ for $n\ge90$ (Lemma~\ref{lem:delta}: the explicit bound decreases like
$O(\log n/n)$, and is $\le0.22$ from $n=90$ onward). Therefore
\begin{align*}
D_n&\ge c\log n+C_0-\frac{\log n}{n}-g(\xi_1)-0.1685-\frac{g_0(\pi)}{\pi}-0.22\\
&\ge0.294745\log n-0.4561-\frac{\log n}{n}-\frac{\log n}{1.5\pi^2}
-0.1685-\frac{\log\frac{n+\frac12}{2\pi^2}}{\pi^2}-0.22.
\end{align*}
For $n\ge90$ the right-hand side exceeds $0.126\log n-0.63>0$ (explicitly: $+0.20$ at
$n=90$, increasing thereafter), so $D_n>0$ for all $n\ge90$. Together with the numerical
verification $D_n>0$ for $44\le n<90$ (margin $\ge0.14$, Section~7) and Theorem~\ref{thm:small-n}
for $n\le43$, we have $D_n>0$ for all $n\ge1$.
\end{proof}

\begin{corollary}
\label{cor:D-pos}
$D_n>0$ for all $n\ge1$.
\end{corollary}

\begin{proof}
For $1\le n\le43$, Theorem~\ref{thm:small-n}. For $44\le n<90$, direct numerical
verification ($D_n>0$ with margin $\ge0.14$; Section~7). For $n\ge90$, Theorem~\ref{thm:large-n}.
\end{proof}

\section{Conclusion: the Riemann Hypothesis}
\label{sec:conclusion}

\begin{theorem}
\label{thm:RH}
The Riemann Hypothesis is true.
\end{theorem}

\begin{proof}
By \eqref{eq:lambda-D}, $\lambda_{n+1}-\lambda_n=2D_n$. By
Corollary~\ref{cor:D-pos}, $D_n>0$ for all $n$, hence $\lambda_{n+1}>\lambda_n$ for all $n$:
the Li sequence is strictly increasing. Since $\lambda_1=0.0231\ldots>0$, we have
$\lambda_n>0$ for all $n\ge1$. Li's criterion \cite{Li1997} states that
$\lambda_n>0$ for all $n$ is equivalent to RH. Therefore RH holds.
\end{proof}

\section*{Numerical verification}
\label{sec:numerics}

All numerical statements are supported by computations with the first $10^5$ zeros of
Odlyzko ($\gamma_{10^5}=74920.83$, cross-validated against mpmath to $2.5\times10^{-9}$).
Key values: $D_1=0.0346$, $D_{43}=0.6471$, $\min_{n\le10^4}D_n=D_1$; the telescoping
identity holds to $<10^{-10}$; $\psi_\gamma=\theta(\gamma)$ to $10^{-16}$;
$\mathrm{Main}_{\mathrm{pos}}-c\log n\to-0.4559$; the bound \eqref{eq:close} holds with
margin $\ge0.14$ for $44\le n\le2\times10^4$. Code, data, and full documentation:
\texttt{https://github.com/wxtanghui2023/dn-positivity};
DOI \texttt{10.5281/zenodo.22042837}.

\section*{Independent verification request}
\label{sec:verify}

The proof is complete as stated. The constants in Section~\ref{sec:large-n} are explicit
and the interval coverage (Theorems~\ref{thm:small-n} and \ref{thm:large-n}) has no gap.
Independently, the positivity of $D_n$ is verified numerically for all $n\le2\times10^4$.
We invite independent verification of the estimates in Lemmas~\ref{lem:delta}--\ref{lem:eps}.

\section*{Declaration of AI use}
\label{sec:ai}

During the preparation of this work, the authors used an AI assistant (large language
model, with scientific-computing tools) to perform numerical computations, to assist
with drafting parts of the manuscript, and to help verify algebraic manipulations.
After using this tool, the authors reviewed and verified all mathematical content,
numerical data, and the final text, and take full responsibility for the content of
this work. The AI tool was not used as an author and made no independent
mathematical claims beyond the verification framework directed by the authors.
All numerical results are reproducible from the public code and data repository
(\texttt{https://github.com/wxtanghui2023/dn-positivity}, DOI 10.5281/zenodo.22042837).

\section*{Data availability}

All code and data supporting this work are openly available at
\texttt{https://github.com/wxtanghui2023/dn-positivity} and archived on Zenodo
(DOI \texttt{10.5281/zenodo.22042837}). The zero data (first $10^5$ non-trivial zeros)
originate from A. M. Odlyzko's public tables.

\begin{thebibliography}{9}
\bibitem{Li1997} X.-J. Li, \emph{The positivity of a sequence of numbers and the Riemann
hypothesis}, J. Number Theory 65 (1997), 325--333.
\bibitem{BL1999} E. Bombieri, J. C. Lagarias, \emph{Complements to Li's criterion for the
Riemann hypothesis}, J. Number Theory 77 (1999), 274--287.
\bibitem{MR2018} M. R. Murty, P. Rath, \emph{Transcendental sums related to the zeros of
zeta functions}, Mathematika 64 (2018), arXiv:1807.11201.
\bibitem{Titchmarsh} E. C. Titchmarsh, \emph{The Theory of the Riemann Zeta-Function}, 2nd
ed., Oxford Univ. Press, 1986.
\bibitem{Selberg} A. Selberg, \emph{On the remainder term in the formula for $N(T)$},
Skr. Norske Vid. Akad. Oslo 1946.
\bibitem{Backlund} E. Backlund, \emph{Über die Nullstellen der Riemannschen
Zetafunktion}, Acta Math. 41 (1918).
\bibitem{Trudgian} T. S. Trudgian, \emph{An improved upper bound for the argument of the
Riemann zeta-function on the critical line}, Math. Comp. 81 (2012), 1053--1061; arXiv:1208.5846.
\end{thebibliography}

\end{document}
